% =====================================================================
%  Research paper template: math-papers
%
%  Usage:  cp templates/research-paper-template.tex NN-project/paper/main.tex
%          cp templates/references-template.bib     NN-project/paper/references.bib
%  Build:  latexmk -pdf -outdir=build main.tex
%
%  Suitable for number theory, combinatorics, algebra and related fields.
%  Replace every <PLACEHOLDER>. Delete sections that do not apply.
% =====================================================================
\documentclass[11pt]{amsart}

% ---------------------------------------------------------------------
% Packages
% ---------------------------------------------------------------------
\usepackage[T1]{fontenc}
\usepackage[utf8]{inputenc}
\usepackage{lmodern}
\usepackage{microtype}
\usepackage[margin=1.1in]{geometry}
\usepackage{amsmath,amssymb,amsthm,mathtools}
\usepackage{booktabs}
\usepackage{enumitem}
\usepackage{xcolor}
\usepackage{graphicx}
\usepackage{listings}          % code listings (optional)
\usepackage[colorlinks=true,linkcolor=black,citecolor=black,urlcolor=blue!50!black]{hyperref}
\usepackage[capitalise,nameinlink]{cleveref}  % must load after hyperref

\lstset{basicstyle=\ttfamily\small,columns=fullflexible,breaklines=true,frame=single}

% ---------------------------------------------------------------------
% Theorem environments
% ---------------------------------------------------------------------
\theoremstyle{plain}
\newtheorem{theorem}{Theorem}[section]
\newtheorem{lemma}[theorem]{Lemma}
\newtheorem{proposition}[theorem]{Proposition}
\newtheorem{corollary}[theorem]{Corollary}
\newtheorem{conjecture}[theorem]{Conjecture}
\newtheorem{problem}[theorem]{Problem}
\newtheorem{claim}[theorem]{Claim}
\newtheorem*{theorem*}{Theorem}

\theoremstyle{definition}
\newtheorem{definition}[theorem]{Definition}
\newtheorem{example}[theorem]{Example}
\newtheorem{observation}[theorem]{Computational observation}

\theoremstyle{remark}
\newtheorem{remark}[theorem]{Remark}
\newtheorem*{notation}{Notation}

% The `proof` environment is provided by amsthm:
%   \begin{proof} ... \end{proof}
%   \begin{proof}[Proof of \cref{thm:main}] ... \end{proof}

\numberwithin{equation}{section}

% ---------------------------------------------------------------------
% Macros
% ---------------------------------------------------------------------
\newcommand{\N}{\mathbb{N}}
\newcommand{\Z}{\mathbb{Z}}
\newcommand{\Q}{\mathbb{Q}}
\newcommand{\R}{\mathbb{R}}
\newcommand{\C}{\mathbb{C}}
\newcommand{\F}{\mathbb{F}}
\DeclareMathOperator{\lcm}{lcm}
\DeclareMathOperator{\rad}{rad}
\newcommand{\abs}[1]{\left\lvert #1 \right\rvert}
\newcommand{\floor}[1]{\left\lfloor #1 \right\rfloor}
\newcommand{\ceil}[1]{\left\lceil #1 \right\rceil}

% Margin note for drafts (comment out the second line for the final version)
\newcommand{\todo}[1]{\textcolor{red}{[TODO: #1]}}
% \renewcommand{\todo}[1]{}

% ---------------------------------------------------------------------
% Front matter
% ---------------------------------------------------------------------
\title[<Short title>]{<Full Title of the Paper>}

\author{<Author Name>}
\address{<Department>, <Institution>, <City>, <Country>}
\email{<email@example.com>}

% \author{<Second Author>}
% \address{...}
% \email{...}

\date{\today}
\keywords{<keyword one>, <keyword two>}
\subjclass[2020]{Primary <11XXX>; Secondary <05XXX>}

\begin{document}

\begin{abstract}
<One paragraph: the problem, its context, what this paper proves (or establishes), and the methods used. State clearly whether the paper contains new results, partial results, computations, or background only.>
\end{abstract}

\maketitle

% Delete this box for a paper with new results that have been checked.
\noindent\fbox{\parbox{\dimexpr\linewidth-2\fboxsep-2\fboxrule}{\small
\textbf{Status.} <Draft / background only / contains new results (checked by author) / independently reviewed / formally verified>.}}

\tableofcontents

% =====================================================================
\section{Introduction}\label{sec:intro}
% =====================================================================
<Context and motivation. Why is the question interesting? Summarize the main results informally, then state the main theorem precisely.>

\begin{theorem}\label{thm:main}
<Main result.>
\end{theorem}

\subsection*{Organization}
<Section-by-section overview.>

% =====================================================================
\section{Mathematical Preliminaries}\label{sec:prelim}
% =====================================================================
\begin{notation}
<Conventions: $\N$, asymptotic notation, logarithms, etc.>
\end{notation}

\begin{definition}\label{def:key}
<Key definition.>
\end{definition}

\begin{lemma}[<Standard result; cite source>]\label{lem:standard}
<Statement.> \cite{<key>}
\end{lemma}

% =====================================================================
\section{Problem Statement}\label{sec:problem}
% =====================================================================
\begin{problem}[<Attribution>]\label{prob:main}
<Exact statement, verified against the original source.>
\end{problem}

<Equivalent formulations, and any conventions (e.g. positive integers) that the statement leaves implicit.>

% =====================================================================
\section{Historical Background}\label{sec:history}
% =====================================================================
<Who posed the problem, when and where (with primary citations). Prizes, with the source and the conditions. Mark unverified details explicitly.>

% =====================================================================
\section{Existing Literature}\label{sec:literature}
% =====================================================================
\subsection{Established results}
<Refereed results, with citations.>

\subsection{Unrefereed reports}
<Preprints, computations and announcements. Label them clearly as unrefereed.>

% =====================================================================
\section{Main Results}\label{sec:results}
% =====================================================================
<Precise statements of all new results. Mark which are new and which are known.>

\begin{proposition}\label{prop:example}
<Statement.>
\end{proposition}

\begin{corollary}\label{cor:example}
<Statement.>
\end{corollary}

% =====================================================================
\section{Proofs}\label{sec:proofs}
% =====================================================================
\begin{proof}[Proof of \cref{prop:example}]
<Complete, rigorous argument.>
\end{proof}

\begin{proof}[Proof of \cref{thm:main}]
<Complete, rigorous argument.>
\end{proof}

% =====================================================================
\section{Computational Experiments}\label{sec:computation}
% =====================================================================
<Algorithms, exact ranges, hardware, run time, software versions, and a link to the code (e.g. ../experiments/). Use exact arithmetic. Never generalize a finite computation into a general claim.>

\begin{observation}\label{obs:example}
<Exactly what was computed, over exactly which range.>
\end{observation}

% =====================================================================
\section{Formal Verification}\label{sec:formal}
% =====================================================================
<Lean 4 / Mathlib version, the exact formal statement (and how it matches \cref{prob:main}), and the axioms used (\texttt{\#print axioms}). Report any \texttt{sorry}, \texttt{native\_decide} or custom axioms.>

\begin{lstlisting}[language={},caption={Formal statement}]
theorem main_result : <statement> := by
  <proof>
\end{lstlisting}

% =====================================================================
\section{Discussion}\label{sec:discussion}
% =====================================================================
<Limitations, why the methods stop where they do, related open questions, and further directions.>

% =====================================================================
\section{Conclusion}\label{sec:conclusion}
% =====================================================================
<Summary of what has and has not been established.>

\section*{Acknowledgements}
<Funding, people, and a disclosure of any AI assistance: which tools were used and for what.>

\bibliographystyle{amsplain}
\bibliography{references}

% \appendix
% \section{<Supplementary material>}

\end{document}
